\section{Tutorial 3: Regularization and Overfitting}

\subsection{Objectives}

This tutorial develops intuition for generalization and regularization:
\begin{itemize}
    \item Understand overfitting and underfitting through examples
    \item Analyze the effect of model complexity
    \item Study the role of regularization
    \item Interpret learning dynamics (early stopping, implicit bias)
\end{itemize}

\medskip

\noindent
\textbf{Focus.}  
Bridge theory and practice through concrete scenarios and reasoning.

\subsection{Exercise 1: Overfitting in Polynomial Regression}

Suppose we fit a polynomial of degree $d$ to noisy data:
\[
y = f(x) + \varepsilon.
\]

\medskip

\noindent
\textbf{Tasks.}
\begin{enumerate}
    \item Describe what happens when $d$ is very small.
    \item Describe what happens when $d$ is very large.
    \item Sketch (qualitatively) training and test error as a function of $d$.
\end{enumerate}

\medskip

\noindent
\textbf{Solution.}

\begin{enumerate}
\item \textbf{Small $d$:}
\begin{itemize}
    \item Model too simple
    \item Cannot capture structure
    \item High bias
\end{itemize}

\item \textbf{Large $d$:}
\begin{itemize}
    \item Model highly flexible
    \item Fits noise
    \item High variance
\end{itemize}

\item \textbf{Behavior:}
\begin{itemize}
    \item Training error decreases monotonically
    \item Test error decreases then increases (U-shape)
\end{itemize}
\end{enumerate}

\medskip

\noindent
\textbf{Key takeaway.}  
Model complexity must be matched to data.

\subsection{Exercise 2: Ridge Regression Effect}

Consider ridge regression:
\[
\min_w \|Xw - y\|^2 + \lambda \|w\|^2.
\]

\medskip

\noindent
\textbf{Tasks.}
\begin{enumerate}
    \item What happens when $\lambda = 0$?
    \item What happens when $\lambda \to \infty$?
    \item Explain how $\lambda$ affects bias and variance.
\end{enumerate}

\medskip

\noindent
\textbf{Solution.}

\begin{enumerate}
\item $\lambda = 0$:
\begin{itemize}
    \item Standard least squares
    \item Potential overfitting
\end{itemize}

\item $\lambda \to \infty$:
\begin{itemize}
    \item $w \to 0$
    \item Model predicts near-constant output
    \item High bias
\end{itemize}

\item Effect:
\begin{itemize}
    \item Increasing $\lambda$ increases bias
    \item Increasing $\lambda$ reduces variance
\end{itemize}
\end{enumerate}

\medskip

\noindent
\textbf{Key takeaway.}  
Regularization controls the bias–variance tradeoff.

\subsection{Exercise 3: Geometry of Regularization}

Consider:
\[
\min_w \|Xw - y\|^2 \quad \text{subject to} \quad \|w\|_2 \leq C.
\]

\medskip

\noindent
\textbf{Tasks.}
\begin{enumerate}
    \item Describe the feasible set.
    \item Explain where the solution lies geometrically.
\end{enumerate}

\medskip

\noindent
\textbf{Solution.}

\begin{enumerate}
\item The feasible set is a ball in $\mathbb{R}^d$.

\item The solution lies where:
\begin{itemize}
    \item level sets of the loss intersect the constraint boundary
\end{itemize}
\end{enumerate}

\medskip

\noindent
\textbf{Insight.}  
Regularization restricts the search space to structured regions.

\subsection{Exercise 4: Implicit Regularization}

Consider overparameterized linear regression:
\[
\min_w \|Xw - y\|^2,
\]
where infinitely many solutions exist.

\medskip

\noindent
\textbf{Tasks.}
\begin{enumerate}
    \item Explain why multiple solutions exist.
    \item Which solution does gradient descent find (starting from $w=0$)?
\end{enumerate}

\medskip

\noindent
\textbf{Solution.}

\begin{enumerate}
\item If $X$ has more parameters than constraints:
\[
Xw = y
\]
has infinitely many solutions.

\item Gradient descent finds the minimum-norm solution:
\[
w^* = \arg\min \|w\| \quad \text{subject to } Xw = y.
\]
\end{enumerate}

\medskip

\noindent
\textbf{Key takeaway.}  
Optimization induces implicit regularization.

\subsection{Exercise 5: Early Stopping}

Suppose we train a model using gradient descent.

\medskip

\noindent
\textbf{Tasks.}
\begin{enumerate}
    \item What happens if we train for too long?
    \item Why can stopping early improve generalization?
\end{enumerate}

\medskip

\noindent
\textbf{Solution.}

\begin{enumerate}
\item Training too long:
\begin{itemize}
    \item Model may fit noise
    \item Overfitting occurs
\end{itemize}

\item Early stopping:
\begin{itemize}
    \item Limits effective complexity
    \item Prevents fitting fine-scale noise
\end{itemize}
\end{enumerate}

\medskip

\noindent
\textbf{Insight.}  
Training time acts as a form of regularization.

\subsection{Exercise 6 (Discussion): Double Descent (Preview)}

\textbf{Question.}  
Classically, test error is U-shaped. Modern deep learning sometimes shows a second descent after interpolation.

\medskip

\noindent
\textbf{Discussion points.}
\begin{itemize}
    \item Overparameterization
    \item Implicit bias of optimization
    \item Role of noise and structure
\end{itemize}

\medskip

\noindent
\textbf{Insight.}  
Generalization behavior in modern models differs from classical intuition.

\subsection{Summary}

\begin{itemize}
    \item Overfitting arises from excessive model flexibility
    \item Regularization controls complexity explicitly and implicitly
    \item Geometry provides insight into constrained optimization
    \item Optimization dynamics influence generalization
\end{itemize}

\medskip

\noindent
\textbf{Preparation for next lectures.}  
We now turn to linear regression, where these ideas can be analyzed exactly and explicitly.